%===============================%
%     IDENTITAS KARYA TULIS     %
%===============================%

% --- Judul ---
%
% Gunakan \judul di bawah ini untuk mengganti judul seperti:
%  - Lembar Jawaban Tutorial Online
%  - Tugas Tutorial Online 1
% Menjadi judul apa pun yang Anda inginkan
%
\newcommand{\judul}{Judul Tulisan Anda yang Disarankan Tidak Terlalu Panjang, Sekitar 10--20 Kata}
\newcommand{\judulENG}{Title of Your Bachelor's Thesis Should Not Be Too Long, Around 10--20 Words}
\newcommand{\judulSebaris}{\begingroup\def\\{ }\judul\endgroup}

% --- Keterangan Maksud atau Tujuan Penulisan Karya Tulis ---
%
% Gunakan \keteranganTujuanKarya di bawah ini untuk menulis 
% pernyataan tujuan dari tulisan yang Anda susun ini.
%
% [Contoh] "Diajukan sebagai salah satu syarat untuk memperoleh
%          gelar Sarjana Sains pada Program Studi ..."
%
\newcommand{\tujuanKarya}{Diajukan sebagai salah satu syarat untuk memperoleh gelar \gelarHarapan.}

% --- Email Korespondensi ---
%
% Umumnya menggunakan NIM@ecampus. Anda dapat
% menggantinya dengan e-mail lain jika dibolehkan.
%
\newcommand{\emailKorespondensi}{\nimMahasiswa @ecampus.ut.ac.id}

% --- Tanggal dan Tahun ---
%
% Gunakan \today untuk tanggal lengkap saat ini.
% Gunakan \the\year untuk tahun ini.
%
\newcommand{\tanggalLengkap}{\today}
\newcommand{\tahun}{\the\year}

% --- Tempat dan Tanggal Lengkap ---
\newcommand{\tempatTanggal}{\daerahMahasiswa, \tanggalLengkap}



%=============================%
%     IDENTITAS MAHASISWA     %
%=============================%===============================%
% Keterangan Fungsi Variabel
% [Kosong]  : Semua
% <-- ENG   : Dokumen Berbahasa Inggris
% <-- LJM   : Lembar Jawaban
% <-- BT    : Buku Tugas
% <-- KTI   : Karya Tulis Ilmiah Sendiri
% <-- KTIDB : Karya Tulis Ilmiah dengan Bimbingan Tutor/Dosen
%=============================================================%

% --- Nama & Nomor Induk ---
\newcommand{\namaMahasiswa}{Yoeru Sandaru}
\newcommand{\nimMahasiswa}{081298765432}

% --- Gelar Akademik yang Ingin Disandang ---
\newcommand{\gelarHarapan}{Sarjana Sains Data} % <-- KTIDB

% --- Program Studi ---
\newcommand{\prodi}{Sains Data}
\newcommand{\prodiENG}{Data Science} % <-- ENG

% --- Fakultas ---
\newcommand{\fakultas}{Sains dan Teknologi}
\newcommand{\fakultasENG}{Science and Technology} % <-- ENG

% --- Daerah ---
\newcommand{\daerahMahasiswa}{Medan}
\newcommand{\daerahKampus}{Medan}
\newcommand{\negaraMahasiswa}{Indonesia} % <-- KTIDB

% --- Perguruan Tinggi ---
\newcommand{\kampus}{Universitas Terbuka}
\newcommand{\kampusENG}{Indonesia Open University} % <-- ENG



%===============================%
%     IDENTITAS TUTOR/DOSEN     %
%===============================%=============================%
% Keterangan Fungsi Variabel
% [Kosong]  : Semua
% <-- ENG   : Dokumen Berbahasa Inggris
% <-- LJM   : Lembar Jawaban
% <-- BT    : Buku Tugas
% <-- KTI   : Karya Tulis Ilmiah Sendiri
% <-- KTIDB : Karya Tulis Ilmiah dengan Bimbingan Tutor/Dosen
%=============================================================%

% --- Nama & Nomor Induk ---
\newcommand{\namaPembimbing}{Revo Wibowo, S.Pd., M.Pd., M.Sc.}
\newcommand{\namaPembimbingPolos}{Revo Wibowo}
\newcommand{\nipPembimbing}{012345678} % <-- KTIDB

% --- Program Studi ---
\newcommand{\prodiPembimbing}{Sains Data} % <-- KTIDB

% --- Fakultas ---
\newcommand{\fakultasPembimbing}{Sains dan Teknologi} % <-- KTIDB

% --- Daerah ---
\newcommand{\daerahPembimbing}{Medan} % <-- KTIDB
\newcommand{\negaraPembimbing}{Indonesia} % <-- KTIDB

% --- Perguruan Tinggi ---
\newcommand{\kampusPembimbing}{Universitas Terbuka} % <-- KTIDB